\begin{textsample}{POS Dim 2 – mg – Score 28.00 – mg/mg\_em\_104.txt}  \label{ex:f2_pos_192}
5.2 A educação \textbf{profissional} em Minas Gerais ( 2000-2020 ) A \textbf{política} de educação \textbf{profissional} em Minas Gerais vem sendo desenhada com avanços e retrocessos pelos programas de governo implementados no âmbito \textbf{federal} e \textbf{estadual} ao longo do tempo.
A consolidação da \textbf{política} é proporcional aos investimentos em estratégias de implementação, infraestrutura, gestão escolar, formação docente, etc.
Segundo Dore e Silva ( 2011 ), nos primeiros anos da década de 1990, a tendência à privatização da educação acentuou -se no Brasil com a retirada do Estado como agente financiador de \textbf{políticas} públicas na área social, atingindo, particularmente, a educação e a cultura.
Minas Gerais seguiu essa tendência da \textbf{política} nacional e, em 1995, o governo mineiro celebrou com o Banco Mundial um contrato de empréstimo para financiamento de projetos para melhoria da \textbf{qualidade} na Educação Básica.
Foram implementados o Programa de Qualidade Total na Educação e o Programa de Melhoria do Ensino Médio, cujo objetivo principal era solucionar o problema de vagas no ensino médio.
Nesse escopo, foi priorizado, na \textbf{rede} pública, o ensino médio com formação propedêutica.
Como consequência, entre 1995-2000, observa -se uma drástica queda em instituições de ensino públicas com \textbf{oferta} de \textbf{cursos} \textbf{técnicos} em Minas Gerais, e um expressivo aumento de matrículas na \textbf{rede} privada ( DORE e SILVA, 2011 ).
Como reflexo dessa \textbf{política}, o ensino \textbf{técnico} profissionalizante iniciou o século XXI com número insignificante de escolas \textbf{estaduais} com \textbf{oferta} de formação \textbf{profissional} na \textbf{rede} pública, principalmente quando comparada à \textbf{rede} particular.
Em 2007, foi criado o Programa de Educação Profissional - PEP por meio do Decreto nº 44.632/2007, com o objetivo de oferecer oportunidades de educação profissionalizante gratuita para alunos do 2º ou 3º ano do Ensino Médio da Rede Pública Estadual ; para aqueles que já concluíram o Ensino Médio e não ingressaram no ensino superior ; e alunos do 1º ou 2º anos dos \textbf{cursos} de Educação de Jovens e Adultos, na modalidade presencial.
No período de 2007 a 2015, o estado de Minas Gerais financiou, por esse programa, o ensino \textbf{técnico} de nível médio para estudantes da \textbf{rede} pública em instituições \textbf{privadas} de ensino em, aproximadamente, 137 \textbf{municípios}, com uma diversidade de mais de 90 \textbf{cursos} \textbf{técnicos} de nível médio.
O Decreto Estadual nº 45.599, de 11 de maio de 2011, \textbf{atualizou} o PEP e criou a Rede Mineira de Formação \textbf{Técnica} de Nível Médio, abrangendo instituições públicas ou privadas, com ou sem fins lucrativos, conveniadas ou credenciadas pela SEE-MG, com \textbf{oferta} de educação \textbf{profissional} \textbf{técnica} de nível médio.
Porém, da mesma forma que as \textbf{políticas} educacionais implantadas em Minas Gerais na década de 1990, o número de escolas com \textbf{oferta} de educação \textbf{profissional} na \textbf{rede} pública de ensino ainda não \textbf{atendeu} a demanda.
Em 2011, o Plano Mineiro de Desenvolvimento Integrado ( PMDI ) 2011-2030 considerou o jovem mineiro como o personagem central para estratégias de desenvolvimento integrado do Estado de Minas Gerais, tendo na educação e, particularmente na educação \textbf{profissional}, os elementos centrais dessa estratégia.
Quatro \textbf{redes} tinham em programas de educação \textbf{profissional} a resposta para diversos problemas do ciclo econômico-social : Rede de Educação e Desenvolvimento Humano, Rede de Desenvolvimento Social e Proteção, Rede de Desenvolvimento Econômico Sustentável, e Rede de Ciência, Tecnologia e Inovação.
Assim, a educação \textbf{profissional} ganhou destaque no PMDI por ser base para o crescimento econômico, para a \textbf{qualificação} para o trabalho e geração de renda.
O resultado foi a criação de um pacto abrangente para o desenvolvimento da Educação Profissional em Minas Gerais que permitisse o alcance das estratégias de Governo previstas no PMDI.
Denominado de Pacto de Desenvolvimento da Educação Profissional de Minas Gerais 2014-2022, foram feitos diagnósticos e estudos sobre desenvolvimento econômico e regional com o objetivo de definir \textbf{diretrizes} e estratégias para orientar a \textbf{política} \textbf{estadual} de educação \textbf{profissional} em Minas Gerais e evitar a desarticulação das ações dentro do Estado.
O Pacto envolveu o governo do Estado, outros entes federados, iniciativa privada, \textbf{entidades} representativas, instituições de ensino e estudantes.
O Pacto foi coordenado pelo Escritório de Prioridades Estratégicas do Governo de Minas Gerais[^26 ], por meio do trabalho conjunto entre servidores públicos e consultoria externa.
Dentro do escopo do programa, consta : - Coleta de informações e projeções relacionadas aos temas da educação \textbf{profissional} e de desenvolvimento econômico em Minas Gerais, com base em relatórios internos e externos, entrevistas, pesquisas etc., e boas práticas nacionais e internacionais. - Engajamento de atores internos ( Governo Estadual ) e externos ( setor produtivo, outros provedores de educação, alunos ) para entendimento de perspectivas e para colaboração para construção do Pacto e comprometimento futuro com as \textbf{diretrizes} nele estabelecidas. - Desenvolvimento de análises para definição da visão e aspiração para educação \textbf{profissional} em Minas Gerais, bem como criação de opções de estrutura de governança e responsabilidades dos atores envolvidos. - Desenvolvimento de um plano de implementação da estratégia e Pacto propostos, com maior granularidade de atividades para 3 setores da economia do estado com potencial para serem os primeiros a implementarem a estratégia desenhada.
Os estudos do Pacto envolveram setores estratégicos do Governo, como o Escritório de Prioridades, as Secretarias de Estado de Educação, Ciência e Tecnologia e Desenvolvimento Social.
Ao final da consultoria, foram entregues 06 produtos para nortear a \textbf{política} de educação \textbf{profissional} em Minas Gerais para o período 2014-2022 : Diagnóstico, Casos de Sucesso, Síntese das Entregas, Estratégia, Governança e Plano de Trabalho.
A partir de 2015, com as mudanças no governo do Estado e na gestão da Secretaria de Estado de Educação, não foram utilizados os resultados do Pacto para as decisões relativas à \textbf{política} de educação \textbf{profissional}.
No cenário nacional, a Lei nº 12.513, de 26 de outubro de 2011, \textbf{instituiu} o Programa Nacional de acesso ao Ensino \textbf{Técnico} e Emprego – Pronatec, com a finalidade de ampliar a \textbf{oferta} de \textbf{cursos} de Educação Profissional e Tecnológica ( EPT ), por meio de programas, projetos e ações de assistência \textbf{técnica} e financeira. [ ^26 ] : Criado em janeiro de 2011, como \textbf{órgão} autônomo e status de Secretaria de Estado e ligado diretamente ao governador, o Escritório de Prioridades Estratégicas tinha por finalidade de contribuir para a definição e a execução das prioridades estratégicas do Governo, assumindo papel colaborador junto aos \textbf{órgãos} e \textbf{entidades} da Administração Pública do Poder Executivo.
Disponível em ( adaptado ).
A \textbf{oferta} de \textbf{cursos} de formação inicial e continuada – FIC e \textbf{cursos} \textbf{técnicos} de nível médio pelo Pronatec é realizada por meio da Bolsa-Formação, regulamentada pela \textbf{Portaria} MEC nº 817, de 13 de agosto de 2015.
O financiamento dos \textbf{cursos} envolve repasse de recursos para pagamento de professores, contratação de apoio administrativo e \textbf{técnico}, assistência estudantil para custeio de despesas de alimentação e transporte do aluno, material didático, etc.
Em 2012, a Secretaria de Estado de Educação de Minas Gerais fez adesão ao Pronatec para a \textbf{oferta} de \textbf{cursos} \textbf{técnicos} de nível médio nas escolas da \textbf{rede} \textbf{estadual}.
Entre 2012 a 2014, a \textbf{rede} \textbf{estadual} de ensino de Minas Gerais \textbf{efetuou} no Sistema Nacional de Informações da Educação Profissional e Tecnológica – SISTEC 42.885 matrículas em \textbf{cursos} \textbf{técnicos} pelo Pronatec.
Em 2017 foi lançado o MedioTec, uma ação dentro do Pronatec para \textbf{ofertar} \textbf{cursos} de educação \textbf{profissional} \textbf{técnica} de nível médio, na forma concomitante, para o aluno das \textbf{redes} públicas \textbf{estaduais} de educação, matriculado no ensino médio regular.
Com novos investimentos para custear \textbf{cursos} \textbf{técnicos} de nível médio, o MedioTec apostava que, após concluir essa etapa de ensino, o estudante estaria apto a se inserir no mundo do trabalho e renda.
Em 2019 e 2020 foi feita a repactuação de saldos Pronatec, envolvendo mais de 1.000 vagas entre \textbf{cursos} \textbf{técnicos} presenciais e vagas para \textbf{cursos} presenciais de Formação Inicial e Continuada – FIC para execução em 2020.
Ela aconteceu no momento em que a SEE reestruturou a \textbf{política} de educação \textbf{profissional} em Minas Gerais.
Partindo de um histórico em que as \textbf{políticas} \textbf{estaduais} para a educação \textbf{profissional} tinham pouca interlocução com as demandas dos setores produtivos, a Secretaria de Estado de Educação, em parceria com a Secretaria de Estado de Desenvolvimento Social de Minas Gerais ( SEDESE ), iniciou, em 2019, estudos para a criação de uma metodologia de prospecção de demanda por \textbf{qualificação} \textbf{profissional} no âmbito do estado.
O estudo tem por objetivo identificar regionalmente, por meio de dados oficiais de vínculos empregatícios, quais as ocupações e setores econômicos com maior probabilidade de contratação de \textbf{profissionais}.
O resultado desse processo é um mapa de demanda que serve de referência para apoiar tecnicamente as decisões gerenciais sobre a \textbf{implantação} de \textbf{cursos} \textbf{técnicos} de nível médio e de \textbf{cursos} de formação inicial e continuada, alterando os critérios para a construção do plano de atendimento da educação \textbf{profissional}.
O mapa de demanda por Educação Profissional em Minas Gerais indica as necessidades de formação \textbf{profissional} por região/município e sugere aquelas mais favoráveis, quanto às expectativas de empregabilidade, para \textbf{implantação} de \textbf{cursos} \textbf{Técnico} de nível médio concomitantes e subsequentes, \textbf{Técnico} integrado ao Ensino Médio, \textbf{cursos} de Formação Inicial e Continuada – FIC e \textbf{cursos} Normal Magistério.
A reestruturação da \textbf{política} de Educação Profissional e Tecnológica ocorre ao mesmo tempo em que é discutida a \textbf{reforma} do Ensino Médio, proposta pela Lei 13.415, de 2017.
As possibilidades da educação \textbf{profissional} se abrem para a flexibilidade dos \textbf{itinerários} \textbf{formativos} e as diferentes configurações do \textbf{currículo} centrado nas escolhas do estudante.

% matched lemmas: atender, atualizar, currículo, curso, diretriz, efetuar, entidade, estadual, federal, formativo, implantação, instituir, itinerário, município, oferta, ofertar, política, portaria, privar, profissional, qualidade, qualificação, rede, reforma, técnico, órgão
\end{textsample}
